\chapter{Endereços}
\label{chapter:addresses}

Um endereço de Monero é uma instrução pública que permite a um remetente criar
uma nova saída para o destinatário. O endereço não é registado directamente na
lista de blocos. A lista de blocos guarda essa saída sob uma chave pública de
utilização única; quem conhecer a chave privada correspondente poderá depois
usá-la como entrada de uma nova transacção. A mesma saída não pode ser gasta
duas vezes.

Por isso, quando dizemos que um endereço tem um saldo, abreviamos uma ideia
mais precisa: uma carteira encontrou saídas que consegue detectar para esse
endereço, determinou os respectivos montantes e ainda não as marcou como
gastas. A conservação dos valores e a ocultação dos montantes serão tratadas
nos capítulos \ref{chapter:pedersen-commitments} e \ref{chapter:transactions};
aqui isolaremos o encaminhamento e a detecção das saídas.

Salvo indicação em contrário, todas as observações sobre a implementação
neste capítulo referem-se à revisão do código Monero identificada no resumo
inicial. A distinção será importante:
algumas formas estão impostas pelo consenso da versão~16, ao passo que outras
são convenções das carteiras para que remetentes e destinatários se entendam.

\section{Chaves}
\label{sec:user-keys}

Bob escolhe dois escalares privados não nulos do corpo escalar,
\(k_B^s,k_B^v\in\mathbb{Z}_l^*\), e publica os pontos
\begin{align*}
 K_B^s&=k_B^sG, & K_B^v&=k_B^vG.
\end{align*}
Como na Secção \ref{ec:keys}, \(G\) gera o subgrupo de ordem prima \(l\).
Chamaremos a \(k_B^s\) \emph{chave privada de gasto} e a \(k_B^v\)
\emph{chave privada de vista}; os expoentes \(s\) e \(v\) indicam
essas duas funções. O endereço principal é o par público
\((K_B^s,K_B^v)\), nunca o par privado.

A separação permite entregar \(k_B^v\), juntamente com o endereço público, a
uma máquina de leitura ou a um auditor sem entregar \(k_B^s\). Uma tal
\emph{carteira só de vista} consegue detectar saídas que lhe foram dirigidas,
reconhecer o sub-endereço usado, descodificar os montantes RingCT e ler um ID
de pagamento curto destinado a Bob. Não consegue calcular a chave privada de
utilização única, criar a respectiva imagem de chave, autorizar um gasto nem
assinar uma transacção. Sem imagens de chave importadas de uma carteira
completa, também não consegue decidir com segurança se uma saída detectada já
foi gasta. A capacidade de observação não confere autoridade de gasto.

As carteiras determinísticas da implementação consultada geram primeiro
\(k^s\) e derivam normalmente \(k^v=\mathcal H_n(k^s)\). É uma convenção de
recuperação da carteira, não uma equação exigida pelo consenso; também é
possível importar dois segredos independentes. Em qualquer caso, os segredos
devem ser escalares canónicos e os pontos públicos esperados pertencem ao
subgrupo gerado por \(G\).

\subsection*{Codificação de um endereço principal}

Para comunicar o par numa representação textual, Monero usa uma variante
própria de Base58 \cite{base-58-encoding,luigi-address}. Na rede principal, a entrada
binária de um endereço principal é
\[
 \operatorname{varint}(18)\mathbin\|K^s\mathbin\|K^v,
\]
pela mesma ordem em que os dois pontos aparecem na estrutura
\texttt{account\_public\_address}. Calcula-se Keccak--256 sobre essa entrada,
juntam-se os primeiros quatro bytes do resultado como soma de verificação e
codifica-se tudo em blocos Base58. O resultado tem 95 caracteres e começa
habitualmente por «4». O número 18, e não esse primeiro carácter isolado, é
o prefixo que identifica simultaneamente a rede e o tipo de endereço.
\marginnote{src/ common/ base58. cpp {\tt encode\_ addr()};
src/ crypto- note\_ basic/ crypto- note\_ basic\_ impl.cpp}

Ao descodificar, a implementação exige Base58 bem formada, soma de verificação
correcta, prefixo da rede e comprimento apropriados, e duas codificações
canónicas de pontos da curva. A soma de quatro bytes detecta enganos de cópia;
não autentica Bob, não cifra o endereço e não prova que alguém conhece as
chaves privadas.

\textbf{Requisitos.} A rotina \texttt{check\_key()} confirma que cada sequência
de 32 bytes representa canonicamente um ponto da curva, mas não testa por si só
se o ponto é não nulo e pertence ao subgrupo de ordem \(l\). Uma carteira
robusta deve exigir estas últimas condições para chaves de endereço, além de
rejeitar escalares nulos ou não canónicos. As chaves produzidas pelo gerador
normal satisfazem-nas por construção.

\section{Endereços ocultos}
\label{sec:one-time-addresses}

Alice não escreve o endereço principal de Bob numa saída. Aplica-lhe antes uma
troca semelhante a Diffie--Hellman e cria uma chave pública nova para cada
saída \cite{cryptoNoteWhitePaper}. O livro chamou historicamente a estes
objectos \emph{endereços ocultos}; diremos também, com maior precisão,
\emph{chaves de saída de utilização única}. Um observador pode ver a chave de
saída, mas não o endereço público do qual ela foi derivada.

\subsection*{Uma saída para um endereço principal}

Isolemos a saída de índice \(t=0\) em que Alice paga \(0{,}123\) XMR a Bob.
Uma transacção RingCT comum da versão~16 tem pelo menos duas saídas, mas a
construção de cada uma pode ser estudada separadamente. Alice conhece
\((K_B^s,K_B^v)\) e procede assim:

\begin{enumerate}
 \item Gera, de modo uniforme, um escalar efémero não nulo
 \(r\in_R\mathbb{Z}_l^*\). A função usada pela implementação rejeita zero e
 produz uma representação escalar canónica.

 \item Calcula a \emph{chave pública de transacção} ordinária
 \[
  R=rG
 \]
 e a derivação partilhada
 \[
  \mathcal D_B=8rK_B^v.
 \]
 O factor 8 é o cofactor da curva. Faz parte da rotina
 \texttt{generate\_key\_derivation()} e não deve ser omitido.

 \item Liga a derivação à posição serializada da saída:
 \begin{align*}
  h_0&=\mathcal H_n\bigl(\mathcal D_B,
             \operatorname{varint}(0)\bigr),\\
  K_0^o&=h_0G+K_B^s.
 \end{align*}
 A função \(\mathcal H_n\) reduz o hash módulo \(l\); o índice é codificado
 como inteiro de comprimento variável, não como o carácter «0».

 \item Coloca \(K_0^o\) no alvo da saída e, segundo a convenção das carteiras,
 coloca \(R\) no campo \texttt{extra}. Na versão~16 o alvo leva ainda a
 etiqueta de vista que definiremos abaixo. O compromisso que oculta o
 montante será explicado na Secção \ref{sec:pedersen_monero}.
\end{enumerate}
\marginnote{src/crypto/ crypto.cpp
{\tt generate\_ key\_ deri- vation()},
{\tt deri- vation\_ to\_ scalar()} e
{\tt derive\_ pu- blic\_ key()}}

Ao percorrer a lista de blocos, Bob lê \(R\) e calcula
\[
 \mathcal D'_B=8k_B^vR.
\]
As duas partes obtêm exactamente os mesmos bytes de derivação, pois
\[
 \mathcal D'_B=8k_B^v(rG)=8r(k_B^vG)=8rK_B^v=\mathcal D_B.
\]
Para o índice \(t=0\), Bob recompõe \(h_0\) e subtrai
\[
 K_0^o-h_0G=K_B^s.
\]
Uma carteira guarda as chaves públicas de gasto que reconhece numa tabela; a
igualdade coloca esta saída na conta principal de Bob. Repetirá o teste para
os índices e chaves públicas de transacção apropriados.

\textbf{Funciona porque} a chave privada correspondente é
\begin{align*}
 k_0^o&=h_0+k_B^s\pmod{l},\\
 K_0^o&=(h_0+k_B^s)G=k_0^oG.
\end{align*}
Alice conhece \(r\), \(\mathcal D_B\), \(h_0\) e \(K_0^o\), mas não conhece
\(k_B^s\) e portanto não obtém \(k_0^o\). Uma carteira só de vista conhece
\(k_B^v\) e reconhece \(K_B^s\), mas continua sem a parcela \(k_B^s\). Só a
carteira completa calcula \(k_0^o\), a imagem de chave e a assinatura
necessárias para gastar. A detecção é uma conclusão local da carteira, não uma
prova pública de propriedade; veremos a autorização do gasto no Capítulo
\ref{chapter:transactions}.

\subsection*{Cofactor, codificações e pontos malformados}

Se as duas chaves públicas provêm de escalares sobre \(G\), todos os pontos
acima já estão no subgrupo de ordem \(l\). A multiplicação por 8 tem uma função
mais geral: elimina a componente de torção de qualquer ponto Edwards
descodificável antes de este ser usado como derivação. Não transforma, porém,
uma chave maliciosa numa chave autenticada. Um ponto de baixa ordem pode
resultar na identidade depois dessa multiplicação, e uma sequência que nem
sequer codifique canonicamente um ponto faz
\texttt{generate\_key\_derivation()} falhar.

As carteiras devem, pois, rejeitar ou saltar chaves públicas de transacção
malformadas e exigir chaves de endereço não nulas no subgrupo primo. Também
devem evitar o caso negligível \(k_t^o=0\), regenerando o segredo efémero se
necessário. A implementação gera \(r\) uniformemente em \(\mathbb{Z}_l^*\), e
os seus pontos honestos têm codificação canónica.

Há aqui uma fronteira de protocolo fácil de perder. O consenso valida que a
chave pública de cada saída é uma codificação de ponto admissível e, na
versão~16, que o alvo tem a variante com etiqueta. Não prova a relação entre
\(K_t^o\), \(R\) e qualquer endereço, nem exige que uma chave pública de
transacção semanticamente válida exista em \texttt{extra}; esses são
contractos de interoperabilidade das carteiras. A rotina de consenso
\texttt{check\_key()} também não impõe, isoladamente, pertença ao subgrupo
primo. Uma transacção deliberadamente defeituosa pode assim ser aceite pelo
consenso e, ainda assim, criar uma saída que a carteira pretendida não
consiga encontrar ou gastar.

\subsection*{Etiquetas de vista}

Testar \(K_t^o-h_tG\) contra uma grande tabela de chaves públicas custa mais do
que comparar um byte. Para cada derivação e índice, o remetente da
implementação consultada calcula
\[
 v_t=\operatorname{primeiroByte}\!\left(
 \mathcal H\bigl(
 \texttt{view\_tag}\mathbin\|\mathcal D_B\mathbin\|
 \operatorname{varint}(t)\bigr)\right),
\]
onde \(\mathcal H\) é Keccak--256 sem redução escalar, e inclui \(v_t\) junto
de \(K_t^o\). O sal ASCII de oito bytes \texttt{view\_tag} separa este uso dos
restantes hashes.
\marginnote{src/crypto/ crypto.cpp
{\tt derive\_ view\_ tag()};
src/crypto- note\_basic/
cryptonote\_ format\_ utils.cpp}

Ao examinar uma saída, Bob calcula primeiro o seu próprio \(v_t\). Se os bytes
diferirem, a carteira salta logo a derivação pública completa. Se coincidirem,
a saída é apenas uma \emph{candidata}: a carteira ainda tem de subtrair
\(h_tG\) e encontrar a chave de gasto resultante na sua tabela. Para saídas
alheias, e sob o comportamento pseudo-aleatório do hash, o filtro elimina em
média 255 de cada 256; cerca de uma em 256 passa como falso positivo e recebe
o teste completo.

Na versão~15 houve um período de transição em que uma transacção podia usar
alvos antigos sem etiqueta ou alvos etiquetados, desde que todas as suas
saídas tivessem o mesmo tipo. Na versão~16, congelada para esta edição, cada
saída tem obrigatoriamente a variante \texttt{txout\_to\_tagged\_key} e uma
etiqueta de exactamente um byte. Saídas históricas sem etiqueta continuam a
ser lidas pelo caminho de varrimento completo.

O consenso verifica a presença do byte, mas não consegue verificar se ele foi
derivado do segredo partilhado. Uma etiqueta errada pode fazer uma carteira
ignorar uma saída que, pela equação de \(K_t^o\), lhe pertenceria. Logo, a
etiqueta é uma optimização de varrimento, não prova de propriedade, não
autorização de gasto e não mecanismo independente de privacidade. Um
observador sem \(r\) nem \(k_B^v\) continua sem poder calcular
\(\mathcal D_B\); quem recebe a chave privada de vista calcula tanto a
etiqueta como o teste completo. Para revelar partes escolhidas do histórico
sem entregar essa chave, recorreremos ao Capítulo
\ref{chapter:tx-knowledge-proofs}.

\subsection{Transacções de múltiplas saídas}
\label{sec:multi_out_transactions}

A transacção corrente costuma pagar ao destinatário e devolver troco ao
remetente. Para transacções RingCT não mineiras, a versão~16 exige por
consenso pelo menos duas saídas. A forma canónica de Bulletproof+ usada nessa
versão agrega no máximo 16; por isso, uma transacção corrente tem
\(2\leq p\leq16\) saídas. As razões de valor e a prova de intervalo ficam para
o capítulo seguinte.

Antes de criar as chaves, a carteira remetente baralha os destinos. O
\emph{índice de saída} \(t\in\{0,\ldots,p-1\}\) é a posição final no vector
serializado \texttt{vout}, e é essa posição que todos os cálculos devem usar.
Para uma transacção cujos destinatários têm endereços principais, um único
segredo fresco \(r\in_R\mathbb{Z}_l^*\) e \(R=rG\) bastam. Para a saída \(t\),
destinada a \((K_t^s,K_t^v)\), definimos
\begin{align*}
 \mathcal D_t&=8rK_t^v=8k_t^vR,\\
 h_t&=\mathcal H_n\bigl(\mathcal D_t,
             \operatorname{varint}(t)\bigr),\\
 K_t^o&=h_tG+K_t^s=(h_t+k_t^s)G,\\
 k_t^o&=h_t+k_t^s\pmod{l}.
\end{align*}
\marginnote{src/crypto- note\_core/
cryptonote\_ tx\_utils.cpp
{\tt construct\_ tx\_ with\_ tx\_ key()}}

Mesmo que Alice pague duas vezes ao mesmo endereço na mesma transacção, os
índices diferentes entram no hash e produzem chaves de saída distintas,
salvo colisão de probabilidade negligível. A etiqueta de vista de cada saída
usa a mesma \(\mathcal D_t\) e o mesmo \(\operatorname{varint}(t)\). O
destinatário recompõe ambos a partir de \(8k_t^vR\); uma carteira completa soma
depois \(k_t^s\) para obter a testemunha privada.

É correcto partilhar \(r\) entre estas saídas da \emph{mesma} transacção. Não
se deve reutilizá-lo noutra transacção: \(R\) repetiria imediatamente, ligando
as duas construções, e repetiria também as derivações para os mesmos
destinatários. Se houver sub-endereços, a escolha de \(R\) e o número de
segredos efémeros mudam; veremos já esse caso na Secção
\ref{sec:subaddresses}.

\section{Sub-endereços}
\label{sec:subaddresses}

Um endereço principal pode originar muitos sub-endereços. Todos conduzem à
mesma carteira, mas Bob pode entregar um diferente a cada cliente, factura ou
serviço e, quando a saída for detectada, saber qual deles foi usado. Ao
contrário de publicar repetidamente o endereço principal, isto evita que um
observador ligue essas identidades apenas por igualdade textual
\cite{MRL-0006-subaddresses}.

Os sub-endereços entraram no software associado à versão~7 do protocolo
\cite{subaddress-pull-request}. A carteira mantém uma tabela das respectivas
chaves públicas de gasto: construir essa tabela requer tempo e memória
proporcionais ao número de sub-endereços, mas cada saída candidata pode ser
classificada por uma consulta à tabela, sem uma nova passagem completa pela
lista de blocos para cada endereço.

\subsection*{Derivação e codificação}

O índice de uma conta e de um sub-endereço é o par
\[
 \iota=(j,i)\in
 \{0,\ldots,2^{32}-1\}\times\{0,\ldots,2^{32}-1\}.
\]
A implementação reserva \(\iota=(0,0)\) para o endereço principal. Para outro
par, codifica \(j\) e \(i\) como inteiros de 32 bits em ordem
\emph{little-endian} e define
\[
 T_{\rm sub}=\texttt{SubAddr}\mathbin\|\texttt{0x00},
\]
uma sequência de oito bytes: os sete caracteres ASCII e o byte NUL final.
A partir das chaves principais de Bob calcula
\begin{align*}
 m_\iota&=\mathcal H_n\bigl(
   T_{\rm sub}\mathbin\|k_B^v\mathbin\|
   \operatorname{LE32}(j)\mathbin\|\operatorname{LE32}(i)\bigr),\\
 k_B^{s,\iota}&=k_B^s+m_\iota\pmod l,\\
 K_B^{s,\iota}&=K_B^s+m_\iota G
                =k_B^{s,\iota}G,\\
 K_B^{v,\iota}&=k_B^vK_B^{s,\iota}.
\end{align*}
O sub-endereço é o par
\((K_B^{s,\iota},K_B^{v,\iota})\).
\marginnote{src/device/
device\_ default.cpp
{\tt get\_ sub- address\_
secret\_ key()} e
{\tt get\_ sub- address()}}

Não se introduz aqui uma nova chave privada
\(k_B^{v,\iota}\) com o papel da chave de vista principal. Uma carteira
completa conhece ambos os factores e sabe que
\(K_B^{v,\iota}=(k_B^v k_B^{s,\iota})G\). Mas não é esse produto que o
algoritmo de varrimento usa: a chave de vista da conta continua a ser
\(k_B^v\), agora aplicada como escalar a \(K_B^{s,\iota}\). Esta distinção é
precisamente o que permite classificar todos os sub-endereços num só
varrimento.

Uma carteira só de vista que possua \(k_B^v\) e \(K_B^s\) consegue recompor
\(m_\iota\), \(K_B^{s,\iota}\) e \(K_B^{v,\iota}\), e preencher a tabela de
detecção. Não consegue obter
\(k_B^{s,\iota}=k_B^s+m_\iota\), porque lhe falta \(k_B^s\).
Por sua vez, quem recebe apenas um sub-endereço público não aprende o endereço
principal. Testar se um sub-endereço e um endereço principal partilham a
mesma chave de vista exigiria distinguir a relação Diffie--Hellman
\((G,K_B^v,K_B^{s,\iota},K_B^{v,\iota})\); a privacidade assenta, portanto,
nas hipóteses de dificuldade usuais do grupo, não em segredo da fórmula.
Divulgar \(k_B^v\) permite fazer o teste e liga todos os sub-endereços da
conta. Reutilizar literalmente o mesmo sub-endereço também o liga por
igualdade.

Na rede principal, a codificação Base58 usa o prefixo numérico 42 seguido de
\(K_B^{s,\iota}\) e \(K_B^{v,\iota}\), e a mesma soma de verificação de quatro
bytes. Tem 95 caracteres e começa habitualmente por «8». Tal como no endereço
principal, o carácter visível não substitui a validação do prefixo, do
comprimento, da soma e dos pontos. Os pontos produzidos honestamente estão no
subgrupo primo; uma carteira deve rejeitar identidades e pontos fora desse
subgrupo. O caso
\(k_B^s+m_\iota=0\pmod l\) tem probabilidade negligível, mas uma implementação
defensiva não deve emitir o endereço degenerado.

\subsection{Enviar para um sub-endereço}

Suponhamos que Alice paga a saída de índice \(t\) ao sub-endereço
\((K_B^{s,\iota},K_B^{v,\iota})\). Escolhe
\(r\in_R\mathbb{Z}_l^*\) e, no caso simples em que esse é o único endereço
distinto entre os destinatários que não são troco, publica
\[
 R=rK_B^{s,\iota}
\]
em vez de \(rG\). A derivação e a chave da saída são
\begin{align*}
 \mathcal D_{B,t}&=8rK_B^{v,\iota},\\
 h_t&=\mathcal H_n\bigl(
       \mathcal D_{B,t},\operatorname{varint}(t)\bigr),\\
 K_t^o&=h_tG+K_B^{s,\iota}.
\end{align*}
A etiqueta de vista é calculada sobre esta mesma
\(\mathcal D_{B,t}\) e o mesmo índice.

Bob usa a chave privada de vista principal:
\begin{align*}
 8k_B^vR
 &=8k_B^v(rK_B^{s,\iota})\\
 &=8r(k_B^vK_B^{s,\iota})
  =8rK_B^{v,\iota}
  =\mathcal D_{B,t}.
\end{align*}
Depois do teste da etiqueta, subtrai \(h_tG\) de \(K_t^o\). O resultado
\(K_B^{s,\iota}\) identifica a linha da tabela e, portanto, o índice
\(\iota\). A carteira completa pode ainda formar
\begin{align*}
 k_t^o&=h_t+k_B^{s,\iota}
       =h_t+k_B^s+m_\iota\pmod l,\\
 K_t^o&=k_t^oG.
\end{align*}
Alice conhece \(r\), as duas chaves públicas e \(h_t\), mas não conhece
\(k_B^{s,\iota}\) e portanto não obtém \(k_t^o\). A carteira só de vista
chega à identificação, mas não a esse segredo.
\marginnote{src/crypto/
crypto.cpp
{\tt derive\_ sub- address\_
public\_ key()};
src/crypto- note\_ core/
crypto- note\_ tx\_ utils.cpp}

\subsubsection*{Chave principal e chaves adicionais}

Uma só fórmula não descreve todos os lotes de destinatários. Na revisão
consultada, a carteira classifica os \emph{endereços distintos que não são de
troco} e escolhe as chaves públicas de transacção assim:

\begin{itemize}
 \item Se todos forem endereços principais, usa uma chave principal
 \(R=rG\), sem chaves adicionais.

 \item Se houver exactamente um sub-endereço distinto e nenhum endereço
 principal entre esses destinatários, usa
 \(R=rK^{s,\iota}\), sem chaves adicionais. Várias saídas para esse mesmo
 sub-endereço podem partilhar a derivação; os seus índices continuam a
 produzir \(h_t\) distintos.

 \item Se misturar endereços principais e sub-endereços, ou se houver mais
 de um sub-endereço distinto, mantém uma chave principal \(R=rG\) e cria um
 segredo fresco \(r_t\in_R\mathbb{Z}_l^*\) para cada posição. A chave adicional
 alinhada com uma saída de endereço principal é \(R_t=r_tG\); para uma saída
 de sub-endereço, é \(R_t=r_tK_t^{s,\iota}\).
\end{itemize}

As chaves adicionais aparecem em \texttt{extra} num vector cuja posição
corresponde ao índice final de \texttt{vout}. A posição reservada ao troco
pode não ser usada para a sua derivação: como o remetente conhece a própria
chave de vista, a implementação deriva o troco da chave principal. Esse
pormenor não autoriza deslocar as restantes posições.

Numa saída que use \(R_t\), o segredo partilhado é
\(8r_tK_t^v=8k_t^vR_t\) para um endereço principal, ou
\(8r_tK_t^{v,\iota}=8k_t^vR_t\) para um sub-endereço. É essa derivação, e não
a da chave principal, que entra em \(h_t\) e na etiqueta de vista da saída.

Ao varrer, a carteira tenta primeiro a derivação da chave principal e, quando
existe, a chave adicional da mesma posição \(t\). Uma chave malformada não
produz uma derivação válida: o caminho de varrimento afectado falha ou salta o
candidato. Quando um auxiliar interno coloca a identidade como valor de
recurso, isso marca a falha e não torna válida a chave original. Cada etiqueta
de vista tem de ser calculada com a derivação que realmente criou a sua saída.
Nem o consenso confirma que o vector está alinhado, nem prova as relações
\(R=rG\), \(R_t=r_tG\) ou \(R_t=r_tK^{s,\iota}\): estas são convenções
necessárias à descoberta pelas carteiras.

\subsubsection*{O limite de Janus}

O teste de recepção acima confirma que a chave de gasto recuperada pertence à
tabela de Bob; não confirma publicamente que as duas metades do endereço
fornecido a Alice foram usadas como um par. Suponhamos que um remetente
malicioso suspeita da relação entre
\((D_1,C_1)=(K_B^{s,1},k_B^vD_1)\) e
\((D_2,C_2)=(K_B^{s,2},k_B^vD_2)\). Pode escolher \(r\) e publicar
\[
 R=rD_2,\qquad
 K_t^o=\mathcal H_n\bigl(8rC_2,\operatorname{varint}(t)\bigr)G+D_1.
\]
Bob calcula \(8k_B^vR=8rC_2\), recupera \(D_1\) e associa a saída ao primeiro
sub-endereço. Se depois revelar ao remetente que recebeu esse pagamento, a
resposta confirma que \(D_1\) e \(D_2\) usam a mesma chave de vista: é o
chamado ataque de Janus \cite{janus-attack}.

Isto não entrega ao atacante a chave privada nem lhe permite gastar a saída;
é uma possibilidade de ligação provocada activamente. Evitar confirmações
externas automáticas reduz o oráculo, mas uma defesa criptográfica completa
exige um contracto adicional entre remetente e carteira receptora. As
propostas incluem publicar material que permita à carteira completa verificar
que a chave de vista e a chave de gasto usadas pelo remetente pertencem ao
mesmo sub-endereço \cite{janus-mitigation-issue-62}. Tal ligação não é hoje
uma regra de consenso, pelo que uma mitigação de carteira deve declarar
exactamente o que verifica e como trata transacções antigas.

\section{Endereços integrados}
\label{sec:integrated-addresses}

Um endereço integrado junta um endereço \emph{principal} a um ID de pagamento
curto escolhido pelo destinatário. Serve, por exemplo, para Bob atribuir um
identificador de oito bytes a uma factura e reconhecer esse identificador
quando Alice paga. Não cria outro conjunto de chaves e não é uma variante de
sub-endereço.

Na rede principal, a entrada Base58 é
\[
 \operatorname{varint}(19)\mathbin\|K_B^s\mathbin\|K_B^v
 \mathbin\|\operatorname{ID}_{8},
\]
seguida da soma de verificação Keccak de quatro bytes. O texto tem 106
caracteres e começa habitualmente por «4». A descodificação exige exactamente
oito bytes de ID. O prefixo integrado só é aceite com o par de um endereço
principal; não existe uma codificação integrada de sub-endereço
\cite{integrated-addresses}.

Estes endereços continuam suportados na revisão consultada por motivos de
compatibilidade. Não devem confundir-se com os antigos IDs de pagamento
autónomos de 32 bytes, escritos em claro em \texttt{extra}: a carteira corrente
recusa criá-los e ignora-os em blocos a partir da versão~12, conservando apenas
caminhos de leitura histórica. Um \texttt{extra} transporta no máximo um ID
curto cifrado, sem indicar a que saída pertence. Por isso, a construção
corrente requer uma única chave de vista de destinatário não troco a quem o
associar; um pagamento agrupado para destinatários distintos não oferece uma
atribuição inequívoca.

\subsubsection*{Codificar}

Se \(r\) é o segredo efémero da chave pública principal
\(R=rG\), Alice forma para o endereço integrado de Bob
\[
 \mathcal D_B=8rK_B^v,
\]
e calcula Keccak--256 sobre os 33 bytes
\(\mathcal D_B\mathbin\|\texttt{0x8d}\). Sejam os primeiros oito bytes
\[
 q=\operatorname{primeiros8}\!\left(
   \mathcal H(\mathcal D_B\mathbin\|\texttt{0x8d})\right).
\]
O valor colocado no nonce de \texttt{extra} é
\[
 c=\operatorname{XOR}(\operatorname{ID}_{8},q).
\]
Aqui não se aplica \(\mathcal H_n\), não há redução módulo \(l\) e não entra
qualquer índice de saída. O byte \texttt{0x8d}, de valor decimal 141, separa
este hash dos hashes das chaves de saída e das etiquetas de vista.
\marginnote{src/device/
device\_ default.cpp
{\tt encrypt\_ payment\_ id()}}

\subsubsection*{Descodificar}

Com a chave pública principal \(R=rG\), Bob recompõe
\[
 \mathcal D'_B=8k_B^vR=\mathcal D_B,\qquad
 \operatorname{ID}_{8}=\operatorname{XOR}\!\left(
 c,\operatorname{primeiros8}\!\left(
   \mathcal H(\mathcal D'_B\mathbin\|\texttt{0x8d})\right)\right).
\]
A igualdade pressupõe, como nas saídas, que \(R\) corresponde ao segredo
efémero usado por Alice. A operação XOR é a da Secção
\ref{sec:XOR_section}.

Esta transformação apenas oculta o ID a quem não conhece a derivação. Não o
autentica, não o liga por consenso a uma saída, não prova que Bob recebeu o
pagamento e não impede Alice de escolher qualquer texto cifrado. A presença do
campo também distingue estruturalmente a transacção. Para reduzir essa
distinção, o construtor da carteira tenta inserir um ID curto nulo cifrado em
certas transacções comuns sem ID explícito, quando há no máximo dois destinos
e uma única chave de vista aplicável. Trata-se de comportamento de carteira,
não de uma regra de consenso; ao descodificá-lo obtêm-se oito bytes nulos.

\section{Endereços de multi-assinatura}
\label{sec:multisignature-addresses}

Uma carteira de multi-assinatura faz várias partes cooperarem para controlar
as chaves descritas neste capítulo. Não existe, no entanto, um prefixo Base58,
um tipo de endereço ou uma espécie de saída de consenso reservados para
«multi-assinatura»: o endereço partilhado contém o par público habitual e o
remetente cria uma chave de saída de utilização única habitual. A diferença
está em como os participantes formam e usam os segredos da carteira.

Na revisão de implementação fixada para este capítulo, o suporte da carteira
é marcado como experimental, vem desactivado e a interface de linha de
comandos exige activá-lo explicitamente com
\texttt{set enable-multisig-experimental 1}; a interface RPC apresenta
igualmente um aviso de perigo. Isto é estado do software, não uma restrição de
consenso. O Capítulo \ref{chapter:multisignatures} desenvolve o protocolo e as
suas hipóteses operacionais.
